% Generated draft: requires mathematical and source review.
\documentclass[11pt]{article}
\usepackage[margin=25mm]{geometry}
\usepackage{fontspec}
\setmainfont{Latin Modern Roman}
\usepackage{amsmath,amssymb,mathtools,unicode-math}
\setmathfont{Latin Modern Math}
\usepackage{enumitem,needspace,xurl,hyperref}
\setlength{\parindent}{0pt}
\setlength{\parskip}{5pt}
\setlength{\emergencystretch}{3em}
\newcommand{\N}{\mathbb{N}}
\newcommand{\Z}{\mathbb{Z}}
\newcommand{\Q}{\mathbb{Q}}
\newcommand{\R}{\mathbb{R}}
\newcommand{\C}{\mathbb{C}}
\newcommand{\F}{\mathbb{F}}
\newcommand{\A}{\mathbb{A}}
\newcommand{\PP}{\mathbb{P}}
\newcommand{\E}{\mathbb{E}}
\newcommand{\eps}{\varepsilon}
\newcommand{\dd}{\,\mathrm{d}}
\begin{document}
\section*{UnsolvedMath ID: 700321}
\textbf{Source problem code:} catalogue entry
\subsection{Definitions and mathematical statement}
For a group \(S\), define its derived series by \(S^{(0)}=S\) and
\[
S^{(i+1)}=[S^{(i)},S^{(i)}]\qquad(i\geq 0),
\]
where \([K,K]\) is the subgroup generated by the commutators \(x^{-1}y^{-1}xy\) with \(x,y\in K\). A group is soluble of derived length at most \(4\) if \(S^{(4)}=\{1\}\). It is finitely generated if \(S=\langle X\rangle\) for some finite subset \(X\subseteq S\). Write \(\operatorname{Aut}(S)\) for the group of automorphisms of \(S\). A faithful action of a group \(H\) on \(S\) by automorphisms is a homomorphism \(\rho:H\to\operatorname{Aut}(S)\) with \(\ker\rho=\{1\}\).

\textbf{Statement (P. M. Neumann; Kourovka Notebook, Problem 5.39).} For every countable group \(H\), including finite groups, there exist a finitely generated group \(S\) and an injective homomorphism
\[
\rho:H\hookrightarrow\operatorname{Aut}(S)
\qquad\text{such that}\qquad S^{(4)}=\{1\}.
\]
The group \(S\) may depend on \(H\).

The supplied registry labels this problem open and also records an explicitly unverified AI proof claim. That claim is not treated here as an established solution; the mathematical statement above is the requested assertion, without a resolution claim.
\subsection{Short English statement}
Every countable group should admit a faithful action by automorphisms on some finitely generated soluble group of derived length at most \(4\). The supplied open label and the accompanying unverified proof claim are source qualifications, not independently verified conclusions.
\subsection{Sources}
\begin{itemize}
\item Local UnsolvedMath registry, permanent ID \(700321\), exact source code \texttt{catalogue entry}, title \emph{Kourovka Notebook Problem 5.39}; supplied complete source record, upstream revision \texttt{d371ee917ac9647fbb06cd6fc529deec6f3dda82}.
\item P. M. Neumann, proposer, 1976. \emph{The Kourovka Notebook (Unsolved Problems in Group Theory)}, 5th issue (1976), Problem 5.39; cited in the supplied record through the 21st issue edition (2026), arXiv:1401.0300v48. The supplied collection description identifies E. I. Khukhro and V. D. Mazurov as editors of Issues 1--21.
\item Background reference supplied in the problem statement: P. M. Neumann, in \emph{Groups--Korea, Pusan, 1988}, Lecture Notes in Mathematics, vol. 1398, Springer, Berlin, 1989, pp. 124--139. The source refers to this contribution for background, particularly the relationship with Problem 8.50; neither the contribution title nor the statement of Problem 8.50 is supplied.
\item Supplied discussion: Achyuth Jayadevan, \emph{Preprint}, 2026, \url{https://zenodo.org/records/22729431}. The record explicitly describes this as an unverified AI solution claim, with the wording that the problem is proved qualified as unverified.
\item Supplied status review dated \texttt{2026-10-16}: canonical status \texttt{open}, evidence category \texttt{unrefereed\_claim}. Its notes report a \(\#\) marker for an unverified AI solution claim in the October 2026 version of the 21st issue, arXiv:1401.0300v48. This is supplied dated curation, not an independently verified assessment of the literature or the claimed proof.
\end{itemize}
\end{document}
